\subsection{A second explicit integral defect: weight twelve at 691}

\begin{corollary}[T-INTEGRAL691]
At $p=691$, every $K_{4k}(\mathbb Z)\{691\}$ vanishes, but the
actual integral tripod Lie algebra is not generated by its odd
Soul\'e classes. With the notation and normalization of
McCallum--Sharifi, one has
\[
 \mathfrak g_{691,12}/\mathfrak h_{691,12}
 \simeq (\mathfrak g_{691}^{\mathrm{ab}})_{12}
 \simeq\mathbb Z/691,
\]
and
\[
 [\sigma_3,\sigma_9]\equiv50[\sigma_5,\sigma_7]
       \pmod{691\mathfrak g_{691,12}}.
\]
Here $\mathfrak h$ denotes the integral Lie subalgebra generated by
the specified odd classes. The displayed coefficient $50$ belongs
to that normalization.
\end{corollary}

\begin{proof}
Vandiver at $691$ lies in the already cited verified range, so
VK-EQUIV gives the $K$-group assertion. Sharifi's cup-product
surjectivity theorem, \cite[Theorem~1.2, p.~2; Corollary~5.9,
p.~27]{Sharifi2007}, gives nontriviality in the irregular component
$(691,12)$. The same paper explicitly makes the application to
Ihara's weight-twelve relation on p.~28.
McCallum--Sharifi \cite[Theorem~9.11, p.~30]{McCallumSharifi2003}
then gives strict inclusion $\mathfrak h_{691,12}\subsetneq
\mathfrak g_{691,12}$ and the displayed congruence. Since $12$ is
the least irregular index, Theorem~9.7 identifies this quotient
with the degree-twelve abelianization. Proposition~9.10 and its
proof, pp.~29--30, identify its nonzero image with the cyclic
group of order $691$ in the stable derivation algebra. This proves
the claimed exact quotient. These are applications of the cited
results, not a new computation of the stable derivation algebra.
\end{proof}

Cup-product surjectivity alone does not assert that the defect for
every irregular prime first occurs at its least irregular index.
Theorem~9.7 additionally requires a nonzero coefficient with
$i+j=m$ and $3\le i<j$; a nonzero pairing whose least positive
representatives sum to $m+p-1$ does not meet that hypothesis.
The two explicit examples $37$ and $691$ have the required
published low-weight computations.
